\section{Conclusion}

We developed \textbf{Camerala}, a privacy-preserving web-based framework for longitudinal affective monitoring, and validated it with a 10-session SCED study. The study revealed a robust **Context $\times$ Appraisal Interaction**: the physiological and behavioral signatures of "Threat" observed in the laboratory were significantly dampened in the home environment.

These findings confirm the "Invariance Breaking" hypothesis: physiological signals do not have fixed meanings. They are dynamically constructed from the interaction of the body, the brain's appraisal, and the environmental context. As we move sensing from the lab to the living room, acknowledging this context-dependency is no longer optional—it is the prerequisite for valid affective computing. SCED provides the rigorous methodological framework needed to map these personalized, context-dependent landscapes.
